\documentclass[../../../main.tex]{subfiles}

\begin{document}

\section{Infinite Descent}

Infinite descent is a powerful proving tool in number theory that
is used very frequently. Many, if not most, problems on Diophantine
equations are proven with infinite descent and it has wide range
of applications from showing that there are no solutions under
a certain condition and finding solutions to advance techniques
like Vieta jumping.

Before we get started with how this technique works, we need to
discuss the foundational principle that infinite descent builds
up on. Consider the following principle.

\begin{theorem}{}{}
The \textbf{well-ordering principle} states that there always
exists a least element in an non-empty set of positive integers.
\end{theorem}

In olympiad, we take the principle above as an axiom and does not
require a proof. Meaning, we can just use ``by well-ordering
principle...'' when we prove by infinite descent.

The well-ordering principle is foundational for infinite descent
because it allows us to have contradiction. Here is what I mean.
When we prove with infinite descent, we assume that there exists
a minimum value in a set. The minimum value could be a number,
sum, and much more depending on the problem. Also, note that we
can assume this because our interests only lies on positive
integers and there cannot be a ``least element'' for rational
and reals. Then, we will infinitely descend, eventually getting
contradiction with our assumption. This is quite hard to understand
without an example, so let's take a look at very famous example
before formally stating the method described.

\begin{problem}{Classic Problem}{}
Show that the equation $x^3 + 2y^3 = 4z^3$ does not have any
solutions in positive integers.
\end{problem}
(\href{https://youtu.be/UJR56wT0NIM}{Video Solution})
\begin{solution}
For the sake of contradiction, assume that there exists positive
integer solutions to the equation and let $(a, b, c)$ be the
solution with the least $a$ by well-ordering principle.
Consider the following congruence.
\[
x^3 + 2y^3 \equiv 4z^3 \equiv 0 \pmod{2}
\]
Therefore, $2 \mid x$. Let $x'$ be an integer such that $x = 2x'$.
Substituting, $8x'^3 + 2y^3 = 4z^3$ and $4x'^3 + y^3 = 2z^3$ are
obtained. By the similar argument, it is evident that $2 \mid y$
and $y = 2y'$ for some integer $y'$. Substituting again,
$4x'^3 + 8y'^3 = 2z^3$ and $2x'^3 + 4y'^3 = z^3$. Therefore,
$2 \mid z$ and there exists some positive integer $z'$ such that
$z = 2z'$. Continuing with the final substitution,
$2x'^3 + 4y'^3 = 8z'^3$ and $x'^3 + 2y'^3 = 4z'^3$ are obtained.
In other words, if $(x, y, z)$ is a solution, $(x', y', z')$ is
also a solution to the equation.

Note that $(a, b, c)$ is assumed to be a solution with the least $a$.
Therefore, $(\frac{a}{2}, \frac{b}{2}, \frac{c}{2})$ is also a
solution to the equation. Because $\frac{a}{2} < a$, there is no
positive integer solutions to the equation.
\end{solution}

As you can see from the example above, we assume the minimum
and show that the infinite descent contradicts with our assumption.

\begin{definition}{}{}
The \textbf{proof by infinite descent} is a technique that shows
contradiction through infinitely descending sequence. Assume that
there exists a solution with minimum condition by well-ordering
principle. Then, construct a strictly and infinitely descending
sequence. The proof completes by contradiction.
\end{definition}

We can expand our idea to see how we can use Pythagorean triples
discussed earlier and infinite descent together to solve problems.

\begin{problem}{}{}
Show that there are no integer triples $(x, y, z)$ where $x \neq 0$
that satisfies $2x^4 + 2x^2y^2 + y^4 = z^2$.

\hfill$\sim$\href{https://math.stackexchange.com/questions/4892263/diophantine-equation-2x4-2x2-y2-y4-z2}{Mathematics Stack Exchange}
\end{problem}
\begin{solution}
First, let $(x, y) = d$. Notice that dividing both sides by $d^4$
gives $\frac{z}{d^2} \in \mathbb{Z}$ and $d^2 \mid z$. Therefore,
let $a = \frac{x}{d}$, $b = \frac{y}{d}$, and $c = \frac{z}{d^2}$
with minimum $a + b$ by well-ordering principle. Substituting to the
original equation divided by $d^2$, $2a^4 + 2a^2b^2 + b^4 = c^2$
and $(a^2 + b^2)^2 + (a^2)^2 = c^2$. Continuing, $m$ and $n$ can be
used to satisfy the Pythagorean triple.

Notice that if $a^2 = m^2 - n^2$, then $a^2 + b^2 = 2mn$. However,
$a^2 + b^2 \equiv 2 \pmod{4}$ or $a^2 + b^2 \equiv 1 \pmod{4}$
since $(a, b) = 1$. In other words, $a^2 = 2mn$, $a^2 + b^2 =
m^2 - n^2$ and $c = m^2 + n^2$ where $(m, n) = 1$.

Continuing, notice that $m^2 = a^2 + b^2 + n^2 \equiv n^2 + 1
\pmod{4}$. If $2 \mid m$, then $m^2 \equiv 2 \pmod{4}$, which is not
possible. Therefore, $2 \mid n$. Let $n = 2p^2$ and $m = q^2$.

Note that $b^2 = m^2 - n^2 - a^2 = (m - n)^2 - 2n^2$ and
$2n^2 = (m - n)^2 - b^2 = 8p^4$. Continuing, let $p = rs$
where $s \leq r$. Rearranging, $8p^4 = (q^2 - 2p^2 + b^2)
(q^2 - 2p^2 - b^2)$ and $2q^2 - 4p^2 = 2r^4 + 4s^2$ or
$2q^2 - 4p^2 = 4r^4 + 2s^2$. Thus, $q^2 - 2r^2s^2 = r^4 + 2s^4$ and
$q^2 = 2s^4 + 2r^2s^2 + r^4$ or $q^2 = 2r^4 + 2r^2s^2 + s^4$.
Notice that this is the same as the original equation where
$(x, y, z) = (s, r, q)$ or $(x, y, z) = (r, s, q)$. Moreover,
because $s + r \leq 2p < 2pq = a < a + b$, there exists no integer
triples $(x, y, z)$ that satisfies the equation.
\end{solution}

Infinite descent have wide range of applications, and we can
extend its application with Vieta's formulas.

\subsection{Vieta Jumping}

Vieta jumping is a technique that has been formally introduced
with the famous problem 1988 IMO Problem 6. Now it is a very
standard and common technique, but back then, it was one of the
revolutionary technique.

The basic idea is to write a given equation as a quadratic
equation with respect to a specific variable, then use the
Vieta's formula to apply infinite descent. To see what I mean,
let's actually solve 1988 IMO Problem 6 together to see how it
works. Also, don't worry! It is often framed as one of the
hardest IMO problems ever, but it is not too difficult when you
understand the solution. Of course the problem is very difficult
because Vieta jumping wasn't really a thing back then, if you
understood the contents in this section, you would be surprised
to see how satisfying this problem is!

\begin{problem}{1988 IMO Problem 6}{1988 IMO Problem 6}
Show that for positive integers $a$ and $b$, if
$\frac{a^2 + b^2}{ab + 1}$ is an integer, then it is a perfect
square.
\end{problem}
(\href{https://youtu.be/rU1e7ttQmh4}{Video Solution})
\begin{solution}
Let $\frac{a^2 + b^2}{ab + 1} = k$ for some integer $k$ and let
$a > b$ without loss of generality. Assume for the sake of
contradiction that $k$ is not the square of an integer.
Rearranging the equation, $a^2 + b^2 = abk + k$ holds and the
following quadratic equation can be written with respect to $a$.
\[
a^2 - bka + b^2 - k = 0
\]
Using Vieta's formulas, the other root $a'$ of the equation satisfy
the following equation if $a$ is a root of the equation.
\[
a'
= bk - a
= \frac{b^2 - k}{a}
\]
Assume by well-ordering principle that $(a, b)$ is a root of the
equation with least $a + b$. Therefore, it suffices to show that
$a' + b < a + b$ and that $a'$ is a positive integer.

First, consider the following inequality for showing that
$a' + b < a + b$.
\[
a'
= \frac{b^2 - k}{a}
< \frac{b^2}{a}
< \frac{b^2}{a}
= a
\]
Therefore, $a' < a$ and $a' + b < a + b$. Continuing, we can show
that $a'$ is always a positive integer. By construction, the following
equation holds.
\[
a'^2 - bka' + b^2 - k = 0
\]
Rearranging the equation, $a'^2 + b^2 = bka' + k$. Assume for the
sake of contradiction that $a' < 0$. Notice that the left-hand
side is positive while the right-hand side $bka' + k = k(ba' + 1) \leq
k(1 - b) \leq 0$. The inequality contradicts with the condition
$bka' + k > 0$ given by the left-hand side. Continuing, if $a' = 0$,
then $b^2 = k$ and this contradicts with the fact that $k$ is not
a perfect square. Thus, $\frac{a^2 + b^2}{ab + 1}$ is always a
perfect square if the value is an integer for positive integers
$a$ and $b$.
\end{solution}

As you can see from the method above, this is pretty simple when
we know what Vieta jumping is! Honestly I am not a huge fan of
the term Vieta jumping because it is infinite descent plus
Vieta's formulas, but I think the term Vieta jumping made it
more intuitive to be used in other type of problems because other
problems can also require jumping from one root to the another
using Vieta's formulas.

This is it for this section! Infinite descent is vary popular
specially in Diophantine equations, and whenever you see
lots of squares in the given equation, it is worth considering
Vieta jumping. In the next section, we will be discussing another
very popular topic in Diophantine equations Pell's equations.

\end{document}


